%% kapa: exam=1 round=30 year=2019 period=1 form=A subject=real_estate subject_ko=부동산학원론 seq=31 number=111
%% source: Q-Net 감정평가사 기출문제 (cst00309) artlSeq=5205205 file=extracted/1008/5205205/2019년 제30회 감정평가사 1차시험 문제지/2019년 제30회 감정평가사 1차시험 문제지/2019년도 제30회 감정평가사 1차 1교시 A형.hwp
%% conversion: kordoc 4.24.51 -> tools/gen1.py (auto); review=proofread-sonnet; flags=none
\begin{question}{111}
민간투자사업의 추진방식에 관한 설명으로 옳지 않은 것은?
\begin{choices}
  \item 사회기반시설의 준공과 동시에 해당 시설의 소유권이 국가 또는 지방자치단체에 귀속되며, 사업시행자에게 일정기간의 시설관리운영권을 인정하는 방식을 BTO방식이라고 한다.
  \item 사회기반시설의 준공과 동시에 해당 시설의 소유권이 국가 또는 지방자치단체에 귀속되며, 사업시행자에게 일정기간의 시설관리운영권을 인정하되, 그 시설을 국가 또는 지방자치단체 등이 협약에서 정한 기간 동안 임차하여 사용ㆍ수익하는 방식을 BTL방식이라고 한다.
  \item 사회기반시설의 준공 후 일정기간 동안 사업시행자에게 해당 시설의 소유권이 인정되며 그 기간이 만료되면 시설소유권이 국가 또는 지방자치단체에 귀속되는 방식을 BOT방식이라고 한다.
  \item BTO방식은 초등학교 교사 신축사업에 적합한 방식이다.
  \item BTL방식은 사업시행자가 최종수요자에게 사용료를 직접 부과하기 어려운 경우 적합한 방식이다.
\end{choices}
\end{question}
\answer{4}
